%% ma: L12-B08-0125
%% lop: 12
%% chuong: 2
%% bai: 8
%% dang: 12.8.24
%% loai: TN4
%% muc: TH
%% dap_an: "D"
%% nguon: "CĐ2.3.pdf giáo viên gửi 27/10/2026 (M07, Lê Minh Khôi – Nguyễn Anh Khôi, Chương II Chủ đề 3 Biểu thức tọa độ của các phép toán vectơ), Câu 40"
%% kien_thuc: "Bài 8 (điều kiện vuông góc)"
%% trang_thai: chinh-thuc
\begin{ex}
Trong không gian $Oxyz$, cho hai vectơ $\vec u=(1;-2;3)$, $\vec v=(2;0;-3)$. Biết rằng vectơ $\vec a=(x;y;4)$ vuông góc với cả hai vectơ $\vec u$, $\vec v$. Tọa độ của vectơ $\vec a$ là
\choice{$\vec a=(-6;9;4)$}{$\vec a=(6;-9;4)$}{$\vec a=(-6;-9;4)$}{\True $\vec a=(6;9;4)$}
\loigiai{$\vec a\cdot\vec u=x-2y+12=0$, $\vec a\cdot\vec v=2x-12=0\Rightarrow x=6$, $y=9$.}
\end{ex}
